\section{A null for task-vector subspace overlap}
\label{sec:method}

\subsection{What is being measured, and against what}

For a pair of task vectors $\Delta W_i, \Delta W_j \in \mathbb{R}^{d_{\mathrm{out}} \times d_{\mathrm{in}}}$
let $V_i^{(k)}$ hold the top-$k$ right singular vectors of $\Delta W_i$, and write
\begin{equation}
\mathcal{O}(i,j) \;=\; \tfrac{1}{k}\,\big\|\, V_i^{(k)\top} V_j^{(k)} \,\big\|_F^2 ,
\label{eq:overlap}
\end{equation}
the mean squared canonical correlation between the two read subspaces.
For two independent uniformly random $k$-subspaces of $\mathbb{R}^{d_{\mathrm{in}}}$ this has
expectation $k/d_{\mathrm{in}}$, which we verify numerically in
Appendix~\ref{app:invariants}; that value is the isotropic null the literature reports against.

The question is not whether $\mathcal{O}$ exceeds $k/d_{\mathrm{in}}$. It does, by an order of
magnitude. We ask how much remains after conditioning on specified marginal structure in each
update. Throughout, \emph{null-reproduced} means reproduced after holding those stated
statistics fixed. It does not identify which process created them: a statistic measured after
fine-tuning can itself contain task-learned structure.

\subsection{The activation-conditioned null}
\label{sec:null}

Let $C_H = \mathbb{E}[h h^\top] = V\operatorname{diag}(w)V^\top$ be the input second moment of the
\emph{frozen base model} at that layer. This is an operational reference geometry, not the claim
that the final update is exactly contained in the frozen-base activation span: activations evolve
during fine-tuning, and optimiser preconditioning can alter the elementary outer-product
gradient geometry.
Retain the first $m_C$ eigenvectors and partition them into blocks $\{\mathcal{B}_b\}$, with
orthogonal projectors $P_b$; the unresolved complement is one additional block. Define
\begin{equation}
 \mathcal{G}=\{Q\in O(d_{\mathrm{in}}):QP_b=P_bQ\ \text{for every }b\}.
 \label{eq:group}
\end{equation}
For each task independently, draw $Q_i$ from Haar measure on this product of orthogonal groups
and set $\widetilde{\Delta W}_i=\Delta W_iQ_i$. This preserves the full singular spectrum because
$\widetilde{\Delta W}_i\widetilde{\Delta W}_i^\top=\Delta W_i\Delta W_i^\top$. If
$\Pi_i=V_i^{(k)}V_i^{(k)\top}$ and $t_{ib}=\operatorname{tr}(P_b\Pi_i)$, then, away from a
singular-value tie at the rank-$k$ boundary, it also preserves every
block occupancy $t_{ib}$ and every block energy $\|\Delta W_iP_b\|_F$. It randomises which
directions within those allocations two tasks select.
It also fixes richer blockwise Gram and principal-angle invariants listed in
Appendix~\ref{app:invariants}, so the attribution is conditional on more than these summaries.

This construction makes the conditional estimand explicit. Independence and Haar invariance give
\begin{equation}
 \mathbb{E}_{\mathcal N}\,\mathcal O(\widetilde{\Delta W}_i,
 \widetilde{\Delta W}_j)
 =\frac{1}{k}\sum_b\frac{t_{ib}t_{jb}}{\operatorname{rank}(P_b)}.
 \label{eq:null-expectation}
\end{equation}
Thus the scalar null-reproduced component is evaluated exactly from the two observed block-occupancy
profiles, not estimated from random draws. For this activation null, sampled matrices are used only when actual null task
vectors are required---for projector reconstruction, positive controls, and external-reference
diagnostics---not for the scalar headline.
For a reported pool, the null-reproduced fraction is
$f_{\mathcal N}=(\overline{\mathcal O}_{\mathcal N}-\overline{k/d_{\rm in}})/
(\overline{\mathcal O}_{\rm raw}-\overline{k/d_{\rm in}})$; averaging over its stated rows
precedes the ratio.

\paragraph{The blocks use a fixed operational coarsening.}
The exact commutant of empirical $C_H$ rotates only its exact multiplicity eigenspaces. Our
$\mathcal G$ is instead the commutant of a resolution-coarsened, truncated block operator: for a
simple retained spectrum its retained action reduces to independent signs, while the deliberately
unresolved complement is one full orthogonal block. Finite calibration data motivate merging
nearby retained directions as well. We estimate eigenvalue errors from eight disjoint folds and
form maximal consecutive runs connected by overlap of adjacent
$\widehat\lambda_j\mathbin{\pm}2\sigma_j$ intervals; Appendix~\ref{app:calibration} gives the
exact split-sample formula.
This is a fixed operational coarsening rule, not a confidence theorem or a unique null. At $1001$
images, across all measured modules, both ViT-B encoders give a median $37$ blocks over $m_C=64$
directions; the complement is
one full block. Appendix~\ref{app:calibration} reports a full $k/m_C/z$ sensitivity grid, the
budget dependence, stored-spectrum limit, and rejected
fixed-rank alternative; the rule, calibration sample, and measured power are all part of the
specification.

\paragraph{Why not a generative null.}
A generative draw with input covariance proportional to $C_H$ instead invents the update--activation
relationship: it concentrates top-$k$ mass toward high-variance activation directions and
overlaps far more than the data in a common-reference diagnostic
(Appendix~\ref{app:generative}). Coordinate permutation fails for the
opposite reason, destroying the occupancy profile despite passing the offset check.

\subsection{The null for low-rank adapters}
\label{sec:lora-null}

For a LoRA update $\Delta W = \tfrac{\alpha}{r} BA$ \citep{hu2022lora} the row space lies inside
$\mathrm{row}(A)$, and equals it whenever $B$ has full column rank. At initialisation, adapters
sharing $A_0$ therefore have the same admissible $r$-dimensional row space. This remains exact
when $A$ is frozen; when $A$ is trained it need not, so we measure its final overlap rather than
assume it. The activation null does not preserve this regime-specific constraint.

Our \emph{factor-conditioned} null keeps each observed $A$ and the singular values of $B$, and
redraws $B\mapsto U\operatorname{diag}(S_B)W^\top$ with independent Haar singular frames. It
preserves $A$ and $B$'s spectrum, but generally not the product spectrum of $BA$; the LoRA
estimand therefore has different invariants from the activation null. The product-spectrum-
preserving alternative over-randomises $A$'s anisotropy, giving $+0.2023$ offset. Appendix
\ref{app:defects} gives the algebra and alternatives. In the selected null's positive control,
the attainable row-space-constrained plant delivers $37.4$--$57.0\%$ of its nominal increment
and the null recovers $100.7$--$101.1\%$ of that delivered increment. This establishes measured
level and response for that intervention, but does not restore product-spectrum preservation.

\subsection{Estimating $C_H$}
\label{sec:cov}

For the vision encoders studied here the exact $d_{\mathrm{in}} \times d_{\mathrm{in}}$ matrix is
affordable and we use it. Appendix~\ref{app:cov} gives a two-pass estimator for the sizes where
it is not, and records why the calibration set is part of the measurement: $C_H$ is an
expectation over some input distribution, so we draw it from a mixture across several of the
benchmark's own datasets and record the composition with every covariance.
